\section{Question 3: Upsampling and Downsampling}

In this section, image resizing techniques, namely Upsampling and Downsampling, have been implemented on the ESP32-CAM.

\subsection{Implementation Details}
The firmware implemented in \texttt{Q3/Q3.ino} provides a web-based interface for real-time image resizing using the ESP32-CAM. The core functionalities include:

\subsubsection{Image Resizing Algorithm}
The \texttt{resizeImage} function implements a \textbf{Nearest Neighbor Interpolation} algorithm to scale images efficiently on the microcontroller:
\begin{itemize}
    \item It calculates \texttt{x\_ratio} and \texttt{y\_ratio} based on the source and target dimensions.
    \item For every pixel in the target image, it maps back to the corresponding coordinate in the source image.
    \item It directly copies the pixel value from the calculated source coordinate without weighted averaging, which minimizes CPU usage and memory bandwidth, operating in $O(W_{new} \times H_{new})$ time complexity.
\end{itemize}

\subsubsection{Web Server Architeture}
The system uses the \texttt{esp\_http\_server} library to host a control interface:
\begin{itemize}
    \item \textbf{Index Handler (root '/'):} Serves an HTML page containing input controls (slider) and an image element. JavaScript is used to dynamically update the image source with the selected scale parameter.
    \item \textbf{Capture Handler ('/capture'):} 
    \begin{enumerate}
        \item Parses the \texttt{scale} query parameter from the URL.
        \item Captures a frame from the camera.
        \item Allocates a new buffer and applies the resizing algorithm.
        \item Generates a BMP header dynamically based on the new dimensions.
        \item Streams the BMP header, color palette (grayscale), and pixel data to the client.
        \item Frees the allocated memory to prevent leakage.
    \end{enumerate}
\end{itemize}

\subsection{Web Interface}
The system publishes the results of its image processing operations via a web server. The ESP32 hosts a web page on the local network, allowing the user to view the processed images instantly through a browser. This allows for visual verification of the algorithm's success.
